CHAPTER 6. CONCLUSIONS

6.1  Summary

This thesis presented LooksmaxAI, a web-based facial harmony analysis system that combines landmark-based anthropometric measurement with feature-based attractiveness analysis. The system was motivated by the limitations of existing face-rating tools, which often provide a final score without enough measurement-level explanation or rely heavily on fixed golden-ratio-style rules. LooksmaxAI addresses this problem by analyzing both front-profile and side-profile images, calculating detailed facial measurements, applying calibrated scoring, and showing results through an interactive dashboard with both geometric harmony assessment and region-specific beauty analysis.

The thesis first reviewed the background of facial anthropometry, cephalometric analysis, landmark detection, 3DDFA_V2, golden-ratio tools, and feature-based attractiveness prediction. This review showed that reliable facial analysis requires both measurable structure and careful interpretation. It also showed that AI-based scoring can be useful, but it should not replace explainable measurement when users need to understand the result.

The implemented system contains a modern full-stack architecture. The frontend uses Next.js, React, TypeScript, Tailwind CSS, Radix UI primitives, and shadcn/ui. The backend uses Next.js API routes and Python child-process execution. The landmark module uses 3DDFA_V2 and FaceBoxes, while the feature-based attractiveness analysis module uses a dual-stream co-attention architecture with MobileNetV2 networks trained on the SCUT-FBP5500 dataset. The system also uses OpenCV, PIL, NumPy, ONNX Runtime, MediaPipe Face Mesh, and NextAuth.js.

The measurement framework covers 62 facial measurements, including 30 front-profile measurements and 32 side-profile measurements. Front-profile analysis includes eyes, brows, nose, mouth, lips, jaw, facial width, facial height, and facial thirds. Side-profile analysis includes nasal architecture, facial convexity, lip position, jaw structure, forehead slope, cranial relationships, and midface projection. These measurements make the final result more transparent than a single attractiveness score.

The HarmonyScore is calculated using a plateau Gaussian scoring model and hierarchical weighted aggregation. Ideal ranges receive full score, deviations are penalized gradually, key metrics receive higher weights, and a critical-flaw penalty prevents severe local imbalance from being hidden by averaging. The system also includes ethnicity and gender-specific calibration so that it does not rely on one universal facial template.

The feature-based attractiveness analysis module operates independently from the HarmonyScore, providing complementary interpretability through region-specific beauty assessment. Using a dual-stream co-attention architecture, the module processes both RGB facial images and 7-channel semantic parsing maps through parallel MobileNetV2-based networks. MediaPipe Face Mesh detects 478 facial landmarks to segment the face into eight anatomical regions (left eye, right eye, nose, mouth, left eyebrow, right eyebrow, skin, hair), and occlusion-based sensitivity analysis quantifies each feature's individual contribution to overall attractiveness.

The module achieves state-of-the-art performance with SRCC of 0.912 and PCC of 0.921 on the SCUT-FBP5500 dataset through 5-fold cross-validation. It outputs both an overall attractiveness score (0-10 scale) and individual feature scores for each region, accompanied by color-coded heatmaps that visualize which facial regions contribute positively (green) or negatively (red) to perceived beauty. With a compact architecture of 7.9 million parameters and inference time of approximately 320 milliseconds per image, the module enables practical real-time applications.

Keeping the feature-based analysis separate from the measurement-based HarmonyScore is one of the main design contributions of the project because it preserves interpretability while still benefiting from learned visual patterns. Users can compare rule-based geometric harmony analysis with learned feature-based attractiveness assessment, and when they disagree, the detailed measurements and feature-specific scores help explain the discrepancy.

Overall, LooksmaxAI demonstrates how clinical-style facial measurement, modern computer vision, deep learning, and web engineering can be integrated into a practical user-facing system. The project provides detailed facial proportion reports, visual measurement overlays, score gauges, strengths and weaknesses, side-profile normalization, interactive landmark adjustment, feature-based attractiveness analysis with region-specific contributions, interpretable heatmaps, and privacy-conscious processing.

6.2  Limitations

Despite the completed implementation, several limitations remain. The first limitation is dependence on image quality. Poor lighting, motion blur, occlusion, facial hair, makeup, hair covering landmarks, non-neutral expression, or strong head rotation can affect landmark detection and measurement accuracy.

The second limitation is dependence on landmark quality. Automatic landmark detection is fast, but it is not perfect. User refinement improves accuracy, but it also requires user attention and may introduce manual placement errors. Future work should evaluate inter-user consistency when different users adjust the same landmarks.

The third limitation is calibration reliability. Ethnicity and gender-specific ideal ranges make the system more flexible, but these ranges should be validated with larger datasets and expert review. Facial variation is complex, and no calibration profile can fully represent every individual.

The fourth limitation is feature-based analysis model bias. The dual-stream co-attention model depends on the SCUT-FBP5500 training dataset containing 5,500 annotated facial images. If the dataset is unbalanced or culturally biased, the model may produce biased predictions. The reported performance metrics (SRCC=0.912, PCC=0.921) are based on 5-fold cross-validation within this dataset, and real-world performance across diverse populations may vary. A complete evaluation should report dataset composition, demographic distribution, rating procedure, and correlation with human judgment across different cultural contexts.

The fifth limitation is feature independence assumption. The occlusion-based sensitivity analysis treats facial features independently, but attractiveness may depend on feature interactions and synergistic effects. For example, eye-eyebrow harmony or the relationship between nose projection and chin projection may not be fully captured by analyzing each region in isolation.

The sixth limitation is static analysis constraints. The system analyzes single images and cannot capture dynamic attractiveness cues such as facial expressions, animations, or video-based assessment. Dynamic features like smile quality, eye contact, and facial movement patterns may significantly influence perceived attractiveness but are not captured by static image analysis.

The seventh limitation is that LooksmaxAI is not a medical tool. The system can support education, cosmetic planning, and personal analysis, but it should not be used for diagnosis or treatment planning without professional supervision. While the measurements are based on established cephalometric and anthropometric principles, the attractiveness assessments are subjective and culturally influenced.

6.3  Suggestion for Future Works

Several improvements can be developed in future versions. The first direction is longitudinal progress tracking. Users could save analysis results over time and compare changes after lifestyle changes, orthodontic treatment, cosmetic procedures, skincare changes, or weight changes. This would require secure storage and clear consent because facial images are sensitive data.

The second direction is surgical and non-surgical simulation. The system could allow users to simulate changes in nose projection, chin projection, jaw angle, lip position, or facial thirds. Such a feature must be designed carefully so that it remains educational rather than misleading. Integration with professional cosmetic consultation services could provide users with realistic expectations and evidence-based recommendations.

The third direction is a recommendation roadmap. Instead of only showing strengths and weaknesses, LooksmaxAI could generate structured improvement suggestions. These suggestions could separate non-surgical options (makeup techniques, hairstyling, skincare routines, facial exercises), grooming or styling options (eyebrow shaping, beard styling, hair color), dental or orthodontic considerations (teeth alignment, bite correction), and professional consultation topics (rhinoplasty, genioplasty, jaw surgery). The feature-based analysis module could prioritize recommendations by identifying which regions have the lowest scores and highest potential for improvement.

The fourth direction is stronger empirical validation. Future work should evaluate landmark accuracy, measurement repeatability, score stability, runtime, user satisfaction, and agreement between feature-based attractiveness scores and human ratings. The evaluation should include diverse images across gender, ethnicity, age, lighting conditions, and pose variations. User studies should assess whether the feature-specific heatmaps and region scores help users understand their facial aesthetics better than holistic scoring alone.

The fifth direction is expanded demographic support. The current system supports several ethnicity categories, but future versions could use more fine-grained calibration or allow adaptive calibration based on population-specific datasets. This would reduce the limitations of broad labels. The feature-based analysis module could be retrained on more diverse datasets to improve generalization across different cultural beauty standards.

The sixth direction is mobile and real-time analysis. A mobile app or webcam-based version could make the workflow more convenient. Real-time analysis would require additional optimization, stronger privacy safeguards, and robust quality checks. The feature-based analysis module's compact architecture (7.9M parameters, ~320ms inference time) makes it suitable for mobile deployment with appropriate optimization.

The seventh direction is better explainability for the feature-based analysis module. Techniques such as saliency visualization, feature attribution, or model uncertainty estimation could help users understand why the AI model produced a particular score for each facial region. Attention visualization could show which specific pixels or patterns within each region most strongly influenced the contribution score. This would make the feature analysis module more transparent and better aligned with the measurement-based philosophy of LooksmaxAI.

The eighth direction is feature interaction modeling. Future versions could develop methods to capture synergistic effects between facial features, such as eye-eyebrow harmony, nose-chin balance, or midface-lower face proportions. This could involve training models that analyze feature combinations rather than treating each region independently, or developing geometric rules that evaluate multi-feature relationships.

The ninth direction is temporal and dynamic analysis. Extending the system to video input would enable analysis of dynamic attractiveness cues such as facial expressions, smile quality, eye contact patterns, and facial movement symmetry. This would require video-based landmark tracking, temporal feature extraction, and models trained on video datasets with dynamic attractiveness annotations.

The tenth direction is personalized aesthetics. Rather than relying solely on population-level averages, the system could adapt to individual user preferences or specific aesthetic goals. Users could provide feedback on which features they want to emphasize or which beauty standards they want to follow, and the system could adjust its recommendations accordingly while still providing objective measurements.
